\scanpage{28}
Sur Topos loc. annelé (commutativement) $X$, on a
\[
K(X,*)=\unsure{\sum|\bigoplus}_{n\geq0}K(B_{\mathfrak S_{nX}})
\]
On a pour $\xi\in K(X)$
\[
\tau^*(\xi)=\sum\tau^n(\xi)\in1+K(X,*)^{\wedge+}=1+\xi+\tau^2\xi\ \text{---}
\]
Pour $\xi\unsure{=|\in}\operatorname{cl}_{K(X)}(E)$, on a
\[
\tau^n(\xi)=\operatorname{cl}_{K(B_{\mathfrak S_{nX}})}\overset{n}{\bigotimes}E
\]
Ceci posé, on a la relation
\[
\tau^n(\xi)=\sum_{\substack{(\alpha_1,\ldots,\alpha_n)\in\mathbf N^n\\1\alpha_1+2\alpha_2+\cdots+n\alpha_n=n}}
\Sigma_{\alpha_1,\ldots,\alpha_n}(\lambda^1\xi,\ldots,\lambda^n\xi)\rho_1^{\alpha_1}\cdots\rho_n^{\alpha_n}
\]
\margin{Si $\xi$ est effectif,
s'agit-il d'un
\unsure{sommant} \unsure{direct}?
(Relation avec le
th. de commutation
\unsure{à la|de la} H. Weyl --- $\unsure{\xi}$?)}

$\rho_i\in K(B_{\mathfrak S_i})=R_{\mathbf Z}(\mathfrak S_i)$ est l'élément
unité de $K(B_{\mathfrak S_i})$, regardé comme él. de
degré $i$ dans
\[
R(1,\unsure{*})=K(e,*)=\sum_{i\geq0}R(\mathfrak S_i)
\]
\diagram{\begin{tikzcd}[ampersand replacement=\&,row sep=small]e\arrow[d]\\\text{topos ponctuel}\end{tikzcd}}
(de sorte que $\rho_1^{\alpha_1}\cdots\rho_n^{\alpha_n}$ est un élément de
$R(1,n)=R(\mathfrak S_n)$), et $\deleted{\unsure{\Sigma_{\alpha_1,\ldots,\alpha_n}}}$
\[
\Sigma_{\alpha_1,\ldots,\alpha_n}\in\mathbf Z[S_1,\text{---},S_n]
\]
est le polynôme \inserted{isobare de poids $n$ ($S_i$ étant de poids $i$)} qui exprime la fonction
symétrique à $n$ variables
\[
\begin{aligned}
&\sum(x_1\text{---}x_{\alpha_1})(x_{\alpha_1+1}\text{---}x_{\alpha_1+\alpha_2})^2\text{---}
\bigl(x_{\unsure{\alpha_1+\cdots+\alpha_{n-1}+1}}\cdots\unsure{x_{\alpha_1+\cdots+\alpha_n}}\unsure{)}^{\unsure{n}}\\
&=\sum x_1^{\unsure{\nu_1}}\text{---}x_n^{\unsure{\nu_n}}\quad(\text{où }\alpha_i=\text{nb des }\nu_j\text{ égaux à }i)
\end{aligned}
\]
comme polynôme en les fonctions sym.
élémentaires $S_1,S_2,\ldots,S_n$ en
\mspage{27}
